\section{Discussion} \label{sec:disc}
Our findings demonstrate that \ac{dohot} can achieve significantly improved
\ac{dns} query performance by tuning Tor circuit configurations,
specifically by reducing the number of hops and by geographically localizing
relays. These optimizations suggest practical pathways to making anonymous
\ac{dns} queries more feasible for everyday use, without fundamentally
compromising privacy guarantees.

\subsection{Performance vs. Anonymity Trade-offs}
The results in Section~\ref{sec:measurement} confirm that three-hop Tor
circuits incur higher latency than two-hop configurations. While Tor's design
conventionally mandates three hops to balance trust and anonymity, our threat
model, focused on \ac{dns} query anonymity, shows that omitting the middle
relay (i.e., using two hops) represent a compelling optimization, particularly
in a latency-sensitive context such as \ac{dns}.

Our comparison between \texttt{torrc} and \texttt{carml+stem} configurations
reveals that default Tor path selection tends to outperform manual circuit
configurations due to bandwidth-weighted relay selection. Nonetheless,
\texttt{carml+stem} offers better control over custom circuit construction
where all relays reside within the same geographic region, significantly
reducing inter-relay latency. 

Reducing the number of hops and specifying the geographic location of relays in
\ac{dohot} does not compromise client privacy within the context of our threat
model for anonymous \ac{dns} queries. However, we do not recommend these
configurations for general Tor usage, as they may weaken anonymity guarantees
in broader applications. Additionally, the anonymity of \ac{dns} queries cannot
be assured in the presence of colluding adversaries under a global threat
model.

The \ac{dotor} approach sacrifices control for simplicity and (expected) minor performance gains. 
Allowing the exit-node to perform the DNS-resolution removes all control thereof. 
Anonymity is instead preserved solely by the Tor-circuit. The main drawbacks here are that the 
query itself might still be logged even though it is anonymized. 
Likewise, should an adversary control either node in a circuit consisting only of a guard 
and exit (e.g. se2se) in addition to the network it could theoretically be possible to correlate
the query and the client.

\subsection{DoHoT in Context of Existing DNS Privacy Protocols}
\ac{dohot} offers a compelling alternative to protocols like \ac{odns} and
\ac{odoh}. Unlike these, it does not require new infrastructure or trusted
third parties, leveraging the mature and widely distributed Tor network
instead. Moreover, since \ac{doh} encrypts \ac{dns} queries end-to-end, exit
relays in \ac{dohot} setups are blind to the query contents, mitigating key
privacy concerns noted in prior studies. 

Schmitt~\emph{et al.}~\cite{ODNS} compared traditional \ac{dns} resolution with
DNS resolution over the Tor network using the Alexa top 10,000 domains. Their
results showed a median resolution time of 31.31 ms for traditional \ac{dns},
compared to 276.76 ms over Tor. It is important to note that their evaluation
used standard Tor configurations, where DNS queries are resolved in plaintext
at the exit relay, unlike \ac{dohot}, which encrypts queries end-to-end using
\ac{doh}. In our measurements, the default three-hop \ac{dohot} configuration
yielded a median resolution time of under 200 ms. When using a two-hop circuit
with both relays near the client, the median dropped below 100 ms.

Schmitt~\emph{et al.} briefly discusses the use of \ac{dns} in Tor in contrast
to \ac{odns}, highlighting several limitations. They note that Tor introduces
significant latency, is susceptible to censorship, and that exit relays may be
held accountable for transmitted content. Additionally, exit relay operators
can observe \ac{dns} queries, and \ac{dns} traffic patterns may be
fingerprinted to compromise client anonymity. Since queries in \ac{dohot} are
encrypted, exit relays can only observe a client communicating with a resolver
without visibility into the actual \ac{dns} content. While Tor censorship
remains a persistent challenge, techniques such as bridge relays provide
potential workarounds. 

Singanamalla~\emph{et al.}~\cite{singanamalla2020} evaluated the performance of
\ac{odoh} in a deployment setting, comparing it with several other protocols
including standard \ac{doh}, proxied \ac{doh}, various \ac{odoh}
configurations, DNSCrypt, and \ac{dohot}. They reported an average resolution
time of approximately 300 ms for \ac{odoh}, while \ac{dohot} averaged 696.40
ms. In contrast, our experiments showed a substantially lower average for
\ac{dohot} at 238.11 ms. With a two-hop circuit using relays located in the
same country as the client, the average resolution time was further reduced to
155.54 ms.

\subsection{Limitations}
However, \ac{dohot} still inherits limitations from Tor, including potential
censorship and reliance on public volunteer infrastructure. Compared to
\ac{odns} or \ac{odoh}, which can offer lower latency and are easier to deploy in
centralized environments, \ac{dohot}'s reliance on Tor's ecosystem might hinder
scalability for very high-throughput scenarios.

Our findings indicate that the middle relay in Tor's conventional three-hop
design may not be necessary for anonymous \ac{dns} resolution within the bounds
of a non-colluding adversary model. However, omitting it reduces the diversity
of the path and may weaken anonymity in settings with more powerful or
coordinated adversaries.

Localizing the entry and exit relays to the same geographic region as the
client significantly reduced \ac{dns} query latency. However, narrowing the
relay selection pool may reduce entropy in path selection, potentially
increasing the risk of relay-based correlation attacks. This strategy should be
applied with caution depending on the user's threat model and required
anonymity level.

Our experiments primarily measured \ac{dns} query latency using NXDOMAIN
lookups. This approach removes caching as a variable and does not fully reflect
real-world usage where cache hits are common.


\iffalse
\begin{table}[h!]
    \centering
    \caption{Comparative Evaluation Dimensions}
    \label{tab:comp}
    \begin{tabular}{|c|c|c|c|}
         \hline
         Approach & Latency & Anonymity & Deployability \\\hline
         \hline
         Do\{T,H,Q\} & Low      & None      & High   \\\hline
         ODNS        & Medium   & Medium    & Low    \\\hline
         ODoH        & Low      & Medium    & Medium \\\hline
         DoHoT 1.0   & High     & High      & High   \\\hline
         DoHoT 2.0   & Medium   & High      & High   \\\hline
    \end{tabular}
\end{table}
\fi


